\section{Spherical caps and projective nets}                        \label{app:caps}

This appendix records the cap estimates used in Sections
\ref{sec:fixed} and \ref{sec:phase}.  They are included to make the dependence
on the two asymptotic regimes explicit.

\subsection{Fixed dimension}

Let $C_d(\eta)=\{u\in\Sph^{d-1}:\arccos|\ip{u}{e_1}|\leq\eta\}$ be a
projective cap.  For $0<\eta<\pi/2$, normalized surface measure gives
\begin{equation}
 \mu_d(C_d(\eta))
 =\frac{2\int_0^\eta\sin^{d-2}\theta\,d\theta}
        {\int_0^\pi\sin^{d-2}\theta\,d\theta}.                    \label{eq:projective-cap}
\end{equation}
For fixed $d$, $\sin\theta\leq\theta$ yields
\[
 \mu_d(C_d(\eta))\leq c_d\eta^{d-1}.
\]
The reverse inequality with another constant holds for
$0<\eta\leq\eta_d$, since $\sin\theta\geq2\theta/\pi$ on
$[0,\pi/2]$.

Take a maximal set of projective points separated by more than $\eta$.
Maximality makes it an $\eta$-net.  Its projective caps of radius $\eta/2$
are disjoint.  The lower cap-volume estimate therefore bounds its cardinality
by $a_d\eta^{-(d-1)}$.  This proves both assertions in
Equation \eqref{eq:caps-fixed}.

\subsection{High-dimensional cap exponent}

For $U$ uniform on $\Sph^{d-1}$, the first coordinate has density
\begin{equation}
 f_d(t)=c_d(1-t^2)^{(d-3)/2},\qquad
 c_d=\frac{\Gamma(d/2)}{\sqrt\pi\,\Gamma((d-1)/2)}.                \label{eq:coord-density}
\end{equation}
Stirling's formula gives $\log c_d=o(d)$.  For fixed $m\in(0,1)$,
\[
 p_d(m)=\int_m^1 f_d(t)\,dt
 \leq c_d(1-m)(1-m^2)^{(d-3)/2}.
\]
This proves
\[
 \limsup_{d\to\infty}\frac1d\log p_d(m)
 \leq\tfrac12\log(1-m^2).
\]
For every fixed $\epsilon\in(0,1-m)$,
\[
 p_d(m)\geq
 \epsilon c_d\big(1-(m+\epsilon)^2\big)^{(d-3)/2}.
\]
First let $d\to\infty$, then let $\epsilon\downarrow0$.  This gives the
matching lower limit and proves Equation \eqref{eq:cap-ldp}.

For completeness, we spell out why the covering theorem supplies the same
exponent.  A cap $\{u:\ip{u}{v}\geq m\}$ has acute angular radius
$\phi=\arccos m$.  It can be represented as the intersection of a sphere of
radius $r=1/\sin\phi>1$ with a unit Euclidean ball.  Theorem 1 of
\citet{dumer2007covering}, with sphere dimension $d-1$, gives a cover whose
density is at most $O(d\log d)$, uniformly over $r>1$.  Dividing this density
by the normalized cap area $p_d(m)$ bounds the number of caps by
\[
 \frac{O(d\log d)}{p_d(m)}=\exp(dI(m)+o(d)).
\]

\section{Details for the nonadaptive lower bounds}                 \label{app:lower}

We verify the simplex normalizations and the move order used in Lemma
\ref{lem:simplex}.

\subsection{Regular-simplex normalization}

Let $z_0,\ldots,z_d\in\R^d$ obey
\[
 \ip{z_i}{z_i}=\frac{d}{d+1},\qquad
 \ip{z_i}{z_j}=-\frac1{d+1}\quad(i\neq j).
\]
These are centered regular-simplex vertices with side length $\sqrt2$ and
circumradius $\sqrt{d/(d+1)}$.  Multiplying each vertex by $\sqrt2\delta$
gives side $2\delta$ and circumradius
\[
 \sqrt2\delta\sqrt{\frac d{d+1}}=J_d\delta.
\]
After rotation and scaling by $1/m$, the circumradius is $J_d\delta/m$.

For any finite vertex set $K$ and direction $v$,
\[
 \width_K(v)=\max_{x,y\in K}\ip{v}{x-y}.
\]
Every difference of two regular-simplex vertices is an edge vector.  Thus the
edge-correlation condition in Lemma \ref{lem:simplex} implies
$\width_{m^{-1}R\Delta}(v)\leq2\delta$.

\subsection{One transcript and one secret}

Fix a deterministic nonadaptive algorithm.  Its complete query set $V$ is a
known function of the algorithm, independent of the secret and answers.  The
rotation in Lemma \ref{lem:simplex}, the scaled simplex $K$, the midpoint
answers, and the algorithm's output on those answers can therefore all be
determined before the interaction begins.  Choose a farthest vertex
\[
 x^\star\in\arg\max_{x\in K}\|x-A(r)\|_2.
\]
The midpoint construction ensures
$|r_v-\ip{v}{x^\star}|\leq\delta$ for every query.  Hence this is a legal
choice of one secret at the start of the game, not a post hoc switch between
secrets.  Since the minimum enclosing ball of a regular simplex is its
circumsphere,
\[
 \|x^\star-A(r)\|_2\geq\rad(K)=\frac{J_d\delta}{m}.
\]

\subsection{Why random rotation is enough}

Let $E$ be the set of $N_d=\binom{d+1}{2}$ unoriented unit edge directions of
a reference simplex.  If $R$ is Haar uniform on the orthogonal group, then for
each fixed $e\in E$, $Re$ is uniform on the sphere.  The different rotated
edges are dependent, but both lower bounds use only the union bound.  In fixed
dimension,
\[
 \Pr\{\rho(Re,v_i)<\eta\text{ for some }e,i\}
 \leq N_dT\mu_d(C_d(\eta)).
\]
In the exponential-dimensional regime,
\[
 \Pr\{|\ip{Re}{v_i}|>m\text{ for some }e,i\}
 \leq2N_dTp_d(m).
\]
Whenever the displayed upper bound is below one, at least one deterministic
rotation has the required simultaneous avoidance property.

\section{Batch refinement details}                                 \label{app:batch}

We give a more explicit derivation of Theorem \ref{thm:batch} from the robust
Jung machinery of \citet{filmus2026optimal}.  Their proof uses noise
$\delta=1/2$, so queried slabs have width one.  Write
\[
 \eta_M=a_dM^{-1/(d-1)},\qquad
 \beta_M=\frac1{\cos\eta_M}-1=O_d(M^{-2/(d-1)}).
\]
After querying the first-batch net, every feasible region $\Phi_0$ has
$\diam(\Phi_0)\leq1+\beta_M$.

The robust Jung theorem in the witness-core form states that, whenever
$\beta_M$ is below a dimension-dependent threshold and the feasible radius has
not already dropped below the Jung benchmark, there is a regular unit simplex
$\Delta_0$ such that
\[
 \operatorname{Witness}(\Phi_0)
 \subseteq\Delta_0+B(C_d\beta_M).
\]
The witness core has the same Chebyshev radius as the feasible region.  Put
$\gamma_0=C_d\beta_M$.

Suppose at the start of a later batch that
\[
 \operatorname{Witness}(\Phi_b)
 \subseteq\Delta_b+B(\gamma_b),
\]
where $\Delta_b$ is a regular unit simplex.  Query all its
$q_d=\binom{d+1}{2}$ unit edge directions in parallel.  The refinement argument
in the proof of Main Theorem 2 of \citet{filmus2026optimal} gives
\[
 \diam(\operatorname{Witness}(\Phi_{b+1}))
 \leq1+4\gamma_b^2.
\]
A second application of robust Jung gives another regular unit simplex and
\[
 \gamma_{b+1}\leq4C_d\gamma_b^2.
\]
For fixed $d,B$, take $M$ large enough that every iterate remains below the
robust-Jung threshold.  Induction yields
\[
 \gamma_{B-1}
 \leq C_{d,B}\beta_M^{2^{B-1}}
 \leq C'_{d,B}M^{-2^B/(d-1)}.
\]
The radius of $\Delta_{B-1}+B(\gamma_{B-1})$ is at most its simplex
circumradius plus $\gamma_{B-1}$.  This proves the normalized result.  General
$\delta$ follows by scaling.

If the radius ever drops below the Jung benchmark during the argument, the
desired upper bound is already true after enlarging $C_{d,B}$.  Thus the stated
case distinction in the robust theorem causes no gap.

\section{Endpoint consequences of the phase curve}                 \label{app:endpoints}

Theorem \ref{thm:phase} is deliberately stated for a fixed
$\alpha\in(0,\infty)$, where all cap estimates are uniform after fixing a small
margin around $m_\alpha$.  Two endpoint consequences follow without a uniform
large-deviation claim.

First, if $\log T/d\to\infty$, fix any $m<1$.  Lemma
\ref{lem:cover-exp} shows that the sphere can eventually be covered using at
most $T$ caps of threshold $m$.  Proposition \ref{prop:upper} gives
\[
 \limsup_{d\to\infty}\frac{\OPT^{\na}_d(T,\delta)}{\delta}
 \leq\frac{\sqrt2}{m}.
\]
Let $m\uparrow1$.  The universal Jung lower limit from
\citet{filmus2026optimal} gives equality with $\sqrt2$.

Second, if $\log T=o(d)$, nonadaptive strategies are a subset of adaptive
strategies.  The subexponential lower bound of \citet{filmus2026optimal}
therefore implies
\[
 \OPT^{\na}_d(T,\delta)-J_d\delta\longrightarrow\infty.
\]
The limiting curve in Theorem \ref{thm:phase} is consistent with both
endpoints.

\section{Computational sanity checks}                               \label{app:checks}

The file \texttt{verify\_geometry.py} performs two independent diagnostics.

First, for $d\in\{2,3,5\}$ it constructs centered regular-simplex coordinates,
draws a fixed nonadaptive query set, searches over seeded random rotations, and
chooses the rotation with the smallest maximum edge-query correlation $m$.  It
then expands the simplex by $1/m$, forms every midpoint response, and checks
\[
 \max_i\width_K(v_i)\leq2\delta,
 \qquad
 \max_{x\in K,i}|\ip{v_i}{x}-r_i|\leq\delta,
 \qquad
 \rad(K)=\frac{J_d\delta}{m}.
\]
All inequalities pass to floating-point tolerance $10^{-10}$.

Second, it evaluates the exact two-sided cap tail using
$U_1^2\sim\operatorname{Beta}(1/2,(d-1)/2)$.  For several values of $\alpha$
and $d\leq500$, it solves for the finite-dimensional correlation whose cap
probability is $e^{-\alpha d}$ and compares it with
$\sqrt{1-e^{-2\alpha}}$.  The absolute discrepancy decreases monotonically in
every tested sequence.  The generated CSV files contain the raw values.

These checks are not invoked in any proof.  They guard the edge-length,
noise-width, and beta-tail normalizations that are easiest to misstate.
